\section{电磁感应}
\sourceproblem{选择1}{36}{42}
\problemtext 两平行长直导线电流等大反向且同步增大，共面圆环位于两线下方，判断感应电流方向。
\workings 取纸面向外为正，按图近线电流向右、远线向左，环内净磁场入纸面且大小增大。楞次定律要求感应磁场出纸面，故电流逆时针。
\crosscheck 环内各点距近线小于距远线，近线贡献占优；电流方向以原图箭头为准。
\sourceproblem{选择2}{36}{42}
\problemtext 通电螺线管右侧有共轴圆环，判断何种操作产生指定方向的感应电流。
\workings 设螺线管对圆环的有向磁通为 $\Phi=M(x)I$，则 $\mathcal E=-M\dot I-I M'(x)\dot x$。接近或增大电流使原磁通大小增加，远离或减小电流使之减少；感应电流的场分别反抗增加或补偿减少。
\crosscheck 原图绕线和圆环箭头需高清确认后才能唯一匹配选项，本条保留方向待核。
\sourceproblem{选择3}{36}{42}
\problemtext 边长 $a$ 的方形线框绕垂直于均匀磁场的轴匀角速 $\omega$ 转动，初始框面与磁场平行，求感应电动势。
\workings 选取使初期磁通增加的法向，$\Phi=Ba^2\sin\omega t$，故 $\mathcal E=-Ba^2\omega\cos\omega t$，大小 $Ba^2\omega|\cos\omega t|$。
\crosscheck 初始磁通为零而电动势大小最大；符号依回路定向。
\sourceproblem{选择4}{36}{42}
\problemtext 导体棒AB在平行于旋转轴的均匀磁场中绕C旋转，$BC=AB/4$，角速度与磁场同向，比较端点电势。
\workings 旋转速度 $\bm v=\bm\omega\times\bm r$，$\bm v\times\bm B=\omega B\bm r_\perp$ 向外。开路稳态 $V(r)-V(C)=B\omega r^2/2$。由于 $AC=3BC$，故 $V_A>V_B$，选A。
\crosscheck 正电荷受磁力向远离轴的方向；不能仅以两端速度方向相反断言等势。
\sourceproblem{选择5}{36}{42}
\problemtext 直导线CD在均匀磁场中平移，求两端动生电动势（图在PDF43）。
\workings 对刚性直线段，$\mathcal E_{CD}=(\bm v\times\bm B)\cdot(\bm r_D-\bm r_C)$。若 $\theta$ 表示速度和磁场夹角，且线段平行于叉积方向，大小为 $Blv|\sin\theta|$；一般还须乘线段与叉积方向夹角的余弦绝对值。
\crosscheck 图中速度、磁场、线段的空间角度关系待高清核定，暂不强配选项。
\sourceproblem{选择6}{37}{43}
\problemtext 两平行导轨上有可滑动导体棒ab、cd，磁场垂直轨面向下，外力使左棒ab向右移动，判断右棒cd趋势。
\workings 左棒右移使闭路面积减小，感应电流产生向下的磁场以维持原磁通。右棒安培力指向右侧，因而cd向右移动，选D。
\crosscheck 感应作用阻碍两棒间距减少；假定导轨与两棒构成导电闭路且无静摩擦锁定。
\sourceproblem{选择7}{37}{43}
\problemtext 单匝线圈形状、尺寸及线性磁介质不变，改变电流，判断自感系数。
\workings 线性介质中自磁通与电流成正比，$L=\Phi/I$ 为几何和介质决定的常数，选D。
\crosscheck 非线性磁介质或磁滞情况下，自感可随工作点改变，题目排除了该情形。
\sourceproblem{选择8}{37}{43}
\problemtext 导体棒ab在水平导轨上向左匀加速，磁场垂直纸面向里，导轨与电容器相接，判断C极板电荷。
\workings 取右、上、纸外为正，$\bm v=-v\bm i,\bm B=-B\bm k$，$\bm v\times\bm B=-vB\bm j$，棒下端为高势端。按图C接上轨，故C带负电；理想导线零电阻下 $|Q|=CBdv$ 随速率增大。
\crosscheck 若回路电阻不可忽略，应解充电微分方程；题设明确忽略导线及棒电阻。
\sourceproblem{选择9}{37}{43}
\problemtext 圆形区域内均匀轴向磁场随时间变化，水平直导线放在图示ab和更远离圆心的a'b'位置，比较感应电动势。
\workings 圆内感生场 $\bm E=-\dot B\,r\bm e_\phi/2$，沿水平弦积分得到有向电动势 $\mathcal E=\dot B\,y\,l/2$，其中 $y$ 是弦的有向高度、$l$ 是导线位于圆内的长度。若两线段长度相等且均在圆内，距中心较远者电动势大小更大。
\crosscheck 若两段是长度不同的完整弦，则还需比较 $|y|l$，不能只比较距圆心的距离；图示端点范围待复核。
\sourceproblem{选择10}{37}{43}
\problemtext 两真空长螺线管长度与单位长度匝数相同，直径比 $d_1/d_2=1/4$，电流相同，求储能比。
\workings $L=\mu_0n^2Sl$，$W=LI^2/2$，故 $W_1/W_2=(d_1/d_2)^2=1/16$，选B。
\crosscheck 同电流下内部场强相同，储能与场区体积成正比。

\sourceproblem{填空1}{37}{43}
\problemtext 无限长密绕螺线管半径 $r$、单位长度匝数 $n$，电流 $i=I_m\cos\omega t$，外部有同轴单匝圆环半径 $R>r$，求感应电动势。
\workings 理想长管内部 $B=\mu_0ni$，外部场忽略。取法向与正电流产生的磁场同向，$\Phi=\mu_0n\pi r^2 I_m\cos\omega t$，故 $\mathcal E=\mu_0n\pi r^2 I_m\omega\sin\omega t$。
\crosscheck 有效磁通面积是管截面而非外环面积。
\sourceproblem{填空2}{37}{43}
\problemtext 图9半圆弧导线abc及直径方向导线绕O在垂直匀强磁场的平面内逆时针转动，求感应电动势和高电势端（续PDF44）。
\workings 取纸外为正，$\bm B=-B\bm k,\bm\omega=\omega\bm k$。刚体各点 $\bm v\times\bm B=-\omega B\bm r$，从任意端点P到Q的动生电动势为
\[
\mathcal E_{PQ}=-\omega B(r_Q^2-r_P^2)/2.
\]
开路稳态中心O比距O为 $R$ 的端点高 $B\omega R^2/2$；同半径两端等势。
\crosscheck 图9细小端点标注待复核；上式不依赖导线弯曲形状，不能将半圆弧长直接代入 $Blv$。
\sourceproblem{填空3}{38}{44}
\problemtext 两相距不远的平面线圈，调整一个线圈的轴线使其与另一线圈轴线垂直，讨论互感。
\workings 若按题设中心连线与另一线圈轴线的对称配置，垂直线圈穿过的净磁通为零，故互感为零。
\crosscheck 单有轴线垂直而无位置对称条件并不足以保证互感为零；本结论依赖题示几何。
\sourceproblem{填空4}{38}{44}
\problemtext 金属圆盘在沿轴向的均匀磁场中以角速度 $\omega$ 旋转，求盘心与半径 $R$ 的边缘间动生电动势。
\workings 当 $\bm\omega$ 与 $\bm B$ 同向时，$\bm v\times\bm B=\omega Br\bm e_r$，故 $V_{\rm edge}-V_O=\int_0^R B\omega r\,\mathrm dr=B\omega R^2/2$，边缘电势较高。
\crosscheck 反转磁场或转向则极性反转；任一从中心到同半径边缘的路径积分相同。
\sourceproblem{填空5}{38}{44}
\problemtext 图11给出分段线性电流随时间的变化，画自感电动势曲线。
\workings $\mathcal E_L=-L\,\mathrm di/\mathrm dt$。图中电流依次下降、上升、下降、上升，因此相应电动势为正、负、正、负的分段常量；每段高度为该段斜率大小乘 $L$。折点处理想导数不定义。
\crosscheck 原图无数值纵轴，不能臆造幅值；电流过零并不意味着自感电动势过零。
\sourceproblem{填空6}{38}{44}
\problemtext 自感0.4 H的线圈通10 A电流，求储存的磁场能量。
\workings $W=LI^2/2=20\,\mathrm J$。
\crosscheck 采用线性自感模型。
\sourceproblem{填空7}{38}{44}
\problemtext 两长螺线管长度及匝数相同，半径比 $r_1:r_2=1:2$，相对磁导率比 $\mu_{r1}:\mu_{r2}=2:1$，通相同电流，求自感及磁能比。
\workings $L_i=\mu_0\mu_{ri}N^2\pi r_i^2/l$，故 $L_1:L_2=1:2$；由 $W_i=L_i I^2/2$，磁能比也为 $1:2$。
\crosscheck 自感随截面积而非半径一次方变化。
\sourceproblem{填空8}{38}{44}
\problemtext 电容 $10.0\,\mu\mathrm F$，板间电压变化率 $2.0\times10^5\,\mathrm{V/s}$，求位移电流。
\workings 理想电容 $I_d=\mathrm dQ/\mathrm dt=C\,\mathrm dU/\mathrm dt$，代入为 $2.0\,\mathrm A$。
\crosscheck 公式适用于电容不变；高清PDF44确认指数为正5。
\sourceproblem{填空9}{38}{44}
\problemtext 根据感生电场的环流定律判断能否像静电场一样引入单值电势概念。
\workings $\oint\bm E_{\rm ind}\cdot\mathrm d\bm l=-\mathrm d\Phi_B/\mathrm dt$ 一般不为零，故一般不能用单值标量电势的负梯度表示整个感生电场。
\crosscheck 特定无旋区域可以局部引入势；结论不否定电磁势中的标量势与矢量势联合表示。
\sourceproblem{计算1}{39}{45}
\problemtext 长直导线与矩形框共面，近边距离 $d$，平行边长 $a$、垂直宽 $b$，电流 $I=I_m\sin\omega t$，求感应电动势。
\workings 取法向与正电流的磁场同向，
\[
\Phi=\frac{\mu_0aI}{2\pi}\ln\frac{d+b}{d},\qquad
|\mathcal E|=\frac{\mu_0a I_m\omega}{2\pi}\ln\frac{d+b}{d}|\cos\omega t|.
\]
\crosscheck 矩形跨越非均匀场，必须积分；固定框不存在额外动生项。
\sourceproblem{计算2}{39}{45}
\problemtext 半径 $r$ 的半圆形导线绕直径在匀强磁场 $B$ 中以角速度 $\omega$ 转动，回路总电阻 $R$，求感应电流。
\workings 若初态半圆平面与磁场垂直且选初始法向同场，活动半圆的磁通为 $B\pi r^2\cos\omega t/2$，固定部分磁通不随时间变，故
\[
i(t)=\frac{B\pi r^2\omega}{2R}\sin\omega t .
\]
\crosscheck 忽略自感，回路定向反转则电流符号反转；若初始转角不同，应增加初相。
\sourceproblem{计算3}{40}{46}
\problemtext 长直导线电流 $I$ 向下，右侧垂直金属杆CD长 $b$，近端距导线 $a$，杆平行导线向上以速率 $v$ 运动，求电动势及高势端。
\workings 取右、上、纸外为正，右侧 $\bm B=\mu_0 I\bm k/(2\pi x)$，故 $\bm v\times\bm B$ 向右。积分得 $V_D-V_C=\mu_0 Iv\ln[(a+b)/a]/(2\pi)$，D端高势。
\crosscheck 若导线电流或运动方向反转，极性反转；对数来自磁场随距离变化。
\sourceproblem{计算4}{40}{46}
\problemtext 折线OEF绕O在垂直匀强磁场平面内逆时针转动，$OE=EF=a,\angle OEF=120^\circ$，求 $U_{OE},U_{OF}$ 并比较三点电势。
\workings 图示磁场入纸面，角速度出纸面。由余弦定理 $OF^2=3a^2$。开路稳态 $V(r)-V_O=-B\omega r^2/2$，故按 $U_{PQ}=V_P-V_Q$ 定义，
\[
U_{OE}=B\omega a^2/2,\qquad U_{OF}=3B\omega a^2/2.
\]
O电势最高，F最低。
\crosscheck 动生场为径向线性场，折线积分只取始末半径；原图箭头方向与有向电势差须一致。
